\section{Conclusion}
\label{sec:conclusion}

The evaluation protocol, not the model, determines what an encrypted-traffic
classification number means. On 5G Traffic we measure a swing of $0.426$
macro-F1 --- from $0.992$ under a stratified random split to $0.567$ under a
session-disjoint one --- with nothing changed but the partition, and the
mechanism is direct: under the standard protocol, a decision tree on the server
address alone scores $0.999$ and one on the capture identifier alone scores
$0.983$.

The correction is not ``stop splitting randomly''. On CESNET-QUIC22 we tested
five protocols at eight seeds each, and random, temporal and
\emph{client}-disjoint splitting are all statistically indistinguishable from
session-disjoint. What costs $0.192$ there is \emph{server}-disjoint splitting.
Four axes null, one decisive, and which one depends on how the corpus was
collected in a way that is invisible from its label space, its schema or its
class balance.

We therefore propose that the axis be chosen by measurement rather than by
convention: fit a single-identifier probe for each identifier a corpus retains,
under each candidate protocol, and adopt the protocol under which they all
collapse toward the majority-class floor. It costs seconds, it is falsifiable
--- if no protocol collapses them, the corpus cannot support the claim being
made of it --- and we release it as a package that runs on any corpus in our
schema.

\subsection*{What the field looks like under this evaluation}

Different from its published account. A 2016 statistical baseline sits within
$0.01$ macro-F1 of a modern neural model on one corpus. On the other, four
independent neural architectures --- convolutional, recurrent,
transformer-based and our own --- fall within $0.05$ of each other while
gradient-boosted trees on thirteen flow statistics sit $0.23$ above all of
them. The identifier that supplies the labels on a widely-used corpus recovers
them at $0.967$, which caps what any result on it can mean.

\subsection*{On reporting what fails}

Five of the six architectural claims we set out to make did not survive their
own pre-registered tests, and we report them with the tests rather than
quietly. Two are worth carrying beyond this paper.

Adversarial nuisance removal removes nothing, at any weight from $0.01$ to
$2.0$, and it took a raw-feature control to see that: on the corpus where the
nuisance is absent the embedding also leaks nothing, and that vacuous result
would have read as near-perfect invariance had we reported it alone.

Distance-based rejection of unknown applications does beat softmax
thresholding, and the metric training we added to produce that advantage is not
what produces it --- ablating every metric objective makes rejection better.
Measuring the embedding geometry explains why: class separation correlates at
$r = -0.88$ with rejection quality. An objective optimising for discrimination
among known classes covers more of the embedding sphere and leaves less empty
space for novelty to fall into. That applies to any metric-learning approach
used for open-set recognition, not only to ours.

A sixth claim was mis-specified rather than false. Enrollment is flat within a
corpus, where prototypes are already correct; across corpora, where zero-shot
transfer fails outright, five labelled flows per class recover 91\% of the
deficit with the encoder frozen.

\subsection*{Limitations and what comes next}

Our central claim rests on two corpora. Two datasets with opposite structure
are suggestive of a general rule and do not establish one, and the honest
statement is that we have shown the axis differs between these two rather than
that it differs in general. The artifact includes a loader for a third,
ISCX-Tor2016, chosen because it is the only public corpus we found offering a
grouping axis that is neither temporal nor per-file: the same activity was
captured simultaneously at a workstation and at a gateway, so \emph{vantage}
becomes a split axis and the comparison becomes three-way. That corpus sits
behind a registration form and we had not obtained it at the time of writing;
the loader is released with every assumption it makes marked and paired with
how to check it.

We release all split manifests, all results including the failed experiments,
the audit that found them, and a two-minute pipeline verification that requires
no data download.
